\documentclass[11pt]{article}
\usepackage[T1]{fontenc}
\usepackage{amsmath,amssymb,amsthm}
\usepackage[margin=1in]{geometry}
\usepackage{booktabs}
\usepackage{hyperref}\hypersetup{pdftitle={Four Boolean Values for Four-Layer Hamilton Starters at r=3 and r=7},pdfauthor={Xiaojian Zeng},pdfsubject={Exact small-parameter Hamilton starter and chain-compatibility certificates},pdfkeywords={Hamilton cycle, directed circulant, balanced starter, chain compatibility, exact enumeration}}

\newtheorem{theorem}{Theorem}
\newtheorem{lemma}[theorem]{Lemma}
\newcommand{\Z}{\mathbb Z}
\newcommand{\Cr}{\mathcal C_4(r)}

\title{Four Boolean Values for Four-Layer Hamilton Starters at
  \(r=3\) and \(r=7\)}
\author{Xiaojian Zeng\\
PBC School of Finance, Tsinghua University, 100190\\
\texttt{xiaojianzeng@foxmail.com}}
\date{}

\begin{document}
\maketitle

\begin{abstract}
For the directed circulant
\(\mathcal C_4(r)=\operatorname{Cay}^{\to}(\Z_{4r},\{1,\ldots,r\})\),
we determine existence and chain compatibility of four-layer balanced
Hamilton starters at the first two parameters congruent to \(3\pmod 4\).
The exact tuple is
\[
  (E_3,C_3,E_7,C_7)=(0,0,1,1).
\]
The negative value at \(r=3\) follows from an exhaustive enumeration of the
\((3!)^4=1296\) balanced successor assignments: only 15 are permutations,
and their cycle types are \(6+6\) or \(4+8\).  At \(r=7\), an explicit
balanced 28-cycle proves existence, and a displayed simple seven-arc path
has all seven development colours.  All computations are exact, and the
positive claims can be checked directly from the integers printed here.
No assertion is made for another value of \(r\) or about uniqueness.
\end{abstract}

\section{Statement and context}

For odd \(r\), identify \(\Z_{4r}\) with the residues modulo \(4r\), and let
\[
  \Cr=\operatorname{Cay}^{\to}(\Z_{4r},[r]),
  \qquad [r]=\{1,\ldots,r\}.
\]
Thus its arcs are \(x\to x+s\), with arithmetic modulo \(4r\).  A directed
Hamilton cycle \(H\) is a \emph{four-layer balanced Hamilton starter} when,
for every \(j\in\Z_4\) and \(s\in[r]\), it contains exactly one step-\(s\)
arc whose tail is congruent to \(j\pmod 4\).  Write \(E_r\) for the assertion
that such an \(H\) exists.

For \(k\in\Z_r\), put \(H_k=H+4k\).  Balance implies that the \(H_k\)
partition all arcs of \(\Cr\): for fixed \((j,s)\), translating the unique
starter arc through \(4\Z_r\) gives exactly all step-\(s\) arcs with tail
residue \(j\).  Following Jiang and Yang~\cite{JiangYang2026}, the starter is
\emph{chain-compatible} when one arc from each \(H_k\) can be selected so
that the selected arcs, in some order, form a simple directed path.  Let
\(C_r\) denote the existence of such a starter and path.  In particular,
\(C_r\Rightarrow E_r\).

Jiang and Yang construct chain-compatible four-layer starters for
\(r\equiv1\pmod4\), \(r\geq9\), and their Section~9 asks whether a different
construction or an obstruction occurs for \(r\equiv3\pmod4\)
\cite{JiangYang2026}.  The present result concerns only \(r=3,7\).  A bounded
search of the literature did not locate these two exact determinations; this
is a search report, not a priority claim.

The finite-certificate organization used below---an exhaustive finite domain,
exact predicates, and a separate replay verifier---has a methodological
analogue in exact finite-library work in discrete geometry
\cite{ZengEtAl2026ExactFinite}.  That paper supplies no graph-theoretic input
to the present proof.

For use with chain compatibility, balance gives a unique
\(\beta_j(s)\in\Z_r\) such that the step-\(s\) starter arc with tail residue
\(j\) has tail \(j+4\beta_j(s)\).  If \(x=j+4t\), then
\begin{equation}\label{eq:colour}
  \chi(x,s)=t-\beta_j(s)\pmod r
\end{equation}
is the unique \(k\) for which \(x\to x+s\) belongs to \(H_k\)
\cite[Lemma~5.2]{JiangYang2026}.  Hence a simple \(r\)-arc path is a
compatible deletion chain exactly when its colours are all distinct.

\begin{theorem}\label{thm:main}
Under the preceding definitions,
\[
  (E_3,C_3,E_7,C_7)=(0,0,1,1).
\]
\end{theorem}

\section{The exhaustive obstruction at \texorpdfstring{\(r=3\)}{r=3}}

\begin{lemma}\label{lem:r3}
There is no four-layer balanced Hamilton starter in \(\mathcal C_4(3)\).
\end{lemma}

\begin{proof}
Let \(s(x)\in\{1,2,3\}\) be the step used at tail \(x\).  For fixed
\(j\in\Z_4\), balance says precisely that
\[
  \bigl(s(j),s(j+4),s(j+8)\bigr)
\]
is a permutation of \((1,2,3)\).  The four residue classes therefore give
all and only \((3!)^4=1296\) balanced assignments.

For every assignment, form the successor function
\[
  f(x)=x+s(x)\pmod {12}.
\]
A Hamilton starter would make \(f\) a permutation consisting of one
12-cycle.  Exhausting the four independent permutations gives the following
exact census.
\begin{center}
\begin{tabular}{lr}
\toprule
outcome for \(f\) & number of assignments\\
\midrule
not bijective & 1281\\
bijection of cycle type \(6+6\) & 12\\
bijection of cycle type \(4+8\) & 3\\
bijection of cycle type \(12\) & 0\\
\bottomrule
\end{tabular}
\end{center}
The ancillary enumeration implements the following literal loop, so there is
no symmetry assumption in the decisive count:
\[
  (\pi_0,\pi_1,\pi_2,\pi_3)\in S_3^4,
  \qquad s(j+4t)=\pi_j(t)\quad(0\leq t<3),
\]
followed by integer evaluation of \(f\), a bijectivity test, and cycle
tracing.  A second standard-library verifier regenerates the unsymmetrized
domain in a separate pass; it obtains the same 1296 assignments, 15
bijections, and cycle-type counts.  The domain description above proves that
no balanced assignment is omitted, and the cycle census contains no
Hamilton cycle.  Thus \(E_3=0\).
\end{proof}

Since chain compatibility includes the existence of a balanced Hamilton
starter, Lemma~\ref{lem:r3} also gives \(C_3=0\).

\section{A starter and compatible chain at \texorpdfstring{\(r=7\)}{r=7}}

Define a step at each tail \(0,1,\ldots,27\), in that order, by
\begin{equation}\label{eq:tailword}
\begin{split}
(&1,6,4,2,5,5,7,1,4,7,1,3,7,2,3,6,\\
 &6,1,2,5,3,4,5,7,2,3,6,4).
\end{split}
\end{equation}
Iterating the resulting successor function from zero gives
\begin{equation}\label{eq:cycle}
\begin{split}
0,&1,7,8,12,19,24,26,4,9,16,22,27,3,5,10,11,14,\\
  &17,18,20,23,2,6,13,15,21,25,0.
\end{split}
\end{equation}
The 28 entries before the final zero are pairwise distinct, so
\eqref{eq:cycle} is a directed Hamilton cycle using only allowed steps.

For a compact balance check, the tails at which steps \(1,\ldots,7\) occur
are, respectively,
\begin{equation}\label{eq:balancetable}
\begin{array}{c|c@{\qquad}c|c}
s&\text{tails}&s&\text{tails}\\ \hline
1&0,7,10,17 & 5&4,5,19,22\\
2&3,13,18,24& 6&1,15,16,26\\
3&11,14,20,25&7&6,9,12,23\\
4&2,8,21,27&&
\end{array}
\end{equation}
Each row contains exactly one tail in every residue class modulo 4.  Thus
\eqref{eq:cycle} is balanced and proves \(E_7=1\).

It remains to check chain compatibility.  The values \(\beta_j(s)\), with
rows indexed by \(j=0,1,2,3\) and columns by \(s=1,\ldots,7\), are
\begin{equation}\label{eq:beta}
\begin{pmatrix}
0&6&5&2&1&4&3\\
4&3&6&5&1&0&2\\
2&4&3&0&5&6&1\\
1&0&2&6&4&3&5
\end{pmatrix}.
\end{equation}
These entries can also be read directly from \eqref{eq:balancetable}: the
tail is \(j+4\beta_j(s)\).  Consider the path
\begin{equation}\label{eq:path}
  0\longrightarrow1\longrightarrow2\longrightarrow3\longrightarrow4
  \longrightarrow5\longrightarrow6\longrightarrow12,
\end{equation}
whose steps are \((1,1,1,1,1,1,6)\).  Its eight vertices are distinct.
Applying \eqref{eq:colour} and \eqref{eq:beta} to the seven arcs gives
\begin{equation}\label{eq:colours}
  (0,3,5,6,1,4,2),
\end{equation}
which is exactly \(\Z_7\).  Hence the path selects one arc from each of
\(H_0,\ldots,H_6\), proving \(C_7=1\).  Together with the two values from
Lemma~\ref{lem:r3}, this proves Theorem~\ref{thm:main}.

\section{Scope and computational disclosure}

The exhaustive claim is only the nonexistence statement at \(r=3\).  At
\(r=7\), the displayed integers prove existence but make no assertion about
the number or uniqueness of starters or compatible chains.  Nothing here
settles another member of the \(r\equiv3\pmod4\) family.

The \(r=3\) census was generated by exact Python integer arithmetic and was
replayed by a separate standard-library verifier.  An OR-Tools CP-SAT model
was used to discover the \(r=7\) witness, but solver feasibility is not used
as evidence: equations \eqref{eq:tailword}--\eqref{eq:colours} provide the
direct certificate.  OpenAI Codex assisted with finite-search design,
implementation, literature screening, adversarial verification, and initial
manuscript preparation.  Language-model output is not mathematical evidence;
the proof rests on the exhaustive finite domain and the exact displayed
certificates.  No proof assistant was used.

\bibliographystyle{plainurl}
\bibliography{references}

\end{document}
